\section*{Abstract}
		Dieses Dokument liefert die vollständige Herleitung energiebasierter Beziehungen für die Stoffmenge (Mol) und die Lichtstärke (Candela) innerhalb des T0-Modell-Frameworks. Entgegen konventioneller Annahmen, dass diese Größen \textit{Nicht-Energie}-Einheiten seien, demonstrieren wir, dass beide strikt aus dem fundamentalen T0-Energieskalierungsparameter $\xipar = 2\sqrt{G} \cdot E$ hergeleitet werden können. Das Mol ergibt sich als dimensionsloses Verhältnis von Gesamtenergieinhalt zu Teilchen-Energieskala, während die Candela auf einen elektromagnetischen Energiefluss ($[E T^{-1}] = [E^2]$ in natürlichen Einheiten) zurückgeführt wird, gewichtet mit dimensionslosen Faktoren und der Konversionskonstante 683~lm/W. Diese Herleitungen etablieren, dass alle 7 SI-Basiseinheiten fundamentale Energiebeziehungen haben und bestätigen Energie als die universelle physikalische Größe, die vom T0-Modell vorhergesagt wird.

	
	
	\section{Einleitung: Das Energie-Universalitätsproblem}
	\label{sec:einleitung}
	
	\subsection{Konventionelle Sicht: \textit{Nicht-Energie}-Einheiten}
	\label{subsec:konventionelle_sicht}
	
	Die Standardphysik kategorisiert SI-Basiseinheiten in solche mit offensichtlichen Energiebeziehungen und solche ohne:
	
	\textbf{Energiebezogene (5/7):} Sekunde, Meter, Kilogramm, Ampere, Kelvin
	\textbf{Nicht-Energie (2/7):} Mol (Teilchenzählung), Candela (physiologisch)
	
	Diese Klassifikation suggeriert fundamentale Grenzen in der Universalität energiebasierter Physik.
	
	\subsection{T0-Modell-Herausforderung}
	\label{subsec:t0_herausforderung}
	
	Das T0-Modell, basierend auf der universellen Energieskalierung:
	\begin{equation}
		\xipar = 2\sqrt{G} \cdot E
		\label{eq:t0_fundamental}
	\end{equation}
	
	sagt vorher, dass \textbf{alle} physikalischen Größen Energiebeziehungen haben sollten. Dieses Dokument löst den scheinbaren Widerspruch auf, indem es energiebasierte Formulierungen für Mol und Candela herleitet.
	
	\section{Fundamentales T0-Energie-Framework}
	\label{sec:t0_framework}
	
	\subsection{Das universelle Zeit-Energie-Feld}
	\label{subsec:universelles_zeit_energie}
	
	Das T0-Modell etabliert, dass alle Physik aus der fundamentalen Beziehung hervorgeht:
	\begin{equation}
		\Tfield = \frac{1}{\max(E(\vec{x},t), \omega)}
		\label{eq:t0_zeitfeld}
	\end{equation}
	
	wobei $E(\vec{x},t)$ die lokale Energieskala und $\omega$ die charakteristische Frequenz repräsentiert.
	
	\subsection{Feldgleichung und Energiedichte}
	\label{subsec:feldgleichung}
	
	Die regierende Feldgleichung in Energieformulierung:
	\begin{equation}
		\nabla^2 \Tfield = -4\pi G \frac{\rhoE(\vec{x},t)}{\EP} \cdot \frac{\Tfield^2}{\tP^2}
		\label{eq:t0_feldgleichung}
	\end{equation}
	
	verbindet Energiedichte $\rhoE(\vec{x},t)$ mit dem Zeitfeld durch universelle Konstanten.
	
	\section{Stoffmenge (Mol): Energiedichte-Ansatz}
	\label{sec:mol_herleitung}
	
	\subsection{Neukonzeption der \textit{Menge}}
	\label{subsec:neukonzeption_menge}
	
	\subsubsection{Traditionelle Teilchenzählung}
	\label{subsubsec:traditionelle_zaehlung}
	
	Konventionelle Definition:
	\begin{equation}
		n_{\text{konventionell}} = \frac{N_{\text{Teilchen}}}{N_A}
		\label{eq:konventionelles_mol}
	\end{equation}
	
	\textbf{Probleme mit diesem Ansatz:}
	\begin{itemize}
		\item Behandelt Teilchen als abstrakte Entitäten
		\item Keine Verbindung zum physikalischen Energieinhalt
		\item Scheinbar dimensionslos
		\item Fehlt fundamentale theoretische Basis
	\end{itemize}
	
	\subsubsection{T0-Modell: Teilchen als Energieanregungen}
	\label{subsubsec:t0_teilchen_energie}
	
	Im T0-Framework sind Teilchen lokalisierte Lösungen der Energiefeldgleichung. Ein \textit{Teilchen} ist charakterisiert durch:
	
	\begin{equation}
		\text{Teilchen} \equiv \text{Lokalisierte Energieanregung mit charakteristischer Skala } \Echar
		\label{eq:t0_teilchen_definition}
	\end{equation}
	
	\subsection{T0-Herleitung der Stoffmenge}
	\label{subsec:t0_mol_herleitung}
	
	\subsubsection{Energieintegrations-Ansatz}
	\label{subsubsec:energieintegration}
	
	Die \textit{Menge} wird zum Verhältnis zwischen Gesamtenergieinhalt und individueller Teilchenenergie:
	
	\begin{equation}
		\boxed{n_{\text{T0}} = \frac{1}{N_A} \int_V \frac{\rhoE(\vec{x},t)}{\Echar} \, d^3x}
		\label{eq:t0_mol_fundamental}
	\end{equation}
	
	\textbf{Physikalische Komponenten:}
	\begin{itemize}
		\item $\rhoE(\vec{x},t)$: Energiedichtefeld aus dem T0-Modell
		\item $\Echar$: Charakteristische Energieskala des Teilchentyps
		\item $V$: Integrationsvolumen, das die Substanz enthält
		\item $N_A$: Ergibt sich aus T0-Energieskalierungsbeziehungen
	\end{itemize}
	
	\subsubsection{Dimensionsanalyse}
	\label{subsubsec:mol_dimensionsanalyse}
	
	\textbf{Dimension:}
	\begin{equation}
		[n_{\text{T0}}] = \frac{[1][\rhoE][L^3]}{[\Echar]} = \frac{[1][E L^{-3}][L^3]}{[E]} = [1]
		\label{eq:mol_wahre_dimension}
	\end{equation}
	
	\textbf{Erklärung:} Ein Verhältnis Gesamtenergie/Einzelenergie ist dimensionslos, und der Normierungsfaktor $N_A$ ist eine reine Zahl; die Stoffmenge ist in dieser Formulierung ein dimensionsloses Energieverhältnis. Eine $[E^2]$-Dimension folgt daraus nicht. Die Zahl $4{,}040\times10^{-33}$ im Abschnitt zur dimensionalen Verifikation ist $N_A \cdot (1\,\text{eV}/E_P)^2 = 6{,}022\times10^{23}/(1{,}221\times10^{28})^2$ und ebenfalls dimensionslos.
	
	\subsection{Verbindung zum T0-Skalierungsparameter}
	\label{subsec:mol_t0_skalierung}
	
	\subsubsection{Energieskala-Beziehung}
	\label{subsubsec:mol_energieskala}
	
	Für Teilchen atomarer Skala:
	\begin{equation}
		\xipar_{\text{atomar}} = 2\sqrt{G} \cdot \Echar \approx 2\sqrt{G} \cdot (1 \text{ eV}) \approx 10^{-28}
		\label{eq:xi_atomar}
	\end{equation}
	
	\subsubsection{Avogadro-Zahl aus T0-Skalierung}
	\label{subsubsec:avogadro_t0}
	
	Als Ansatz:
	\begin{equation}
		N_A^{(\text{T0})} = \left(\frac{\Echar}{\EP}\right)^{-2} \cdot \mathcal{C}_{\text{T0}}
		\label{eq:avogadro_t0_vorhersage}
	\end{equation}
	
	wobei $\mathcal{C}_{\text{T0}}$ eine dimensionslose Konstante aus der T0-Feldgeometrie sein soll. Die Beziehung ist dimensionslos und damit mit dem dimensionslosen Mol verträglich; $\mathcal{C}_{\text{T0}}$ ist jedoch nicht hergeleitet, die Beziehung bleibt ein Ansatz.
	
	\section{Lichtstärke (Candela): Energiefluss-Wahrnehmung}
	\label{sec:candela_herleitung}
	
	\subsection{Neukonzeption der \textit{Lichtstärke}}
	\label{subsec:neukonzeption_lichtstaerke}
	
	\subsubsection{Traditionelle physiologische Definition}
	\label{subsubsec:traditionelle_lichtstaerke}
	
	Konventionelle Definition:
	\begin{equation}
		I_{\text{konventionell}} = \frac{\Phi_v}{\Omega}, \qquad \Phi_v = 683 \text{ lm/W} \times \Phi_{\text{radiometrisch}} \times V(\lambda)
		\label{eq:konventionelle_candela}
	\end{equation}
	
	wobei $V(\lambda)$ die Augenempfindlichkeitsfunktion des Menschen ist, $\Phi_v$ der Lichtstrom in lm und $\Omega$ der Raumwinkel (Lichtstärke in cd $=$ lm/sr).
	
	\textbf{Probleme mit diesem Ansatz:}
	\begin{itemize}
		\item Abhängig von menschlicher Physiologie
		\item Keine fundamentale physikalische Basis
		\item Willkürliche Normierung (683 lm/W)
		\item Begrenzt auf schmalen Wellenlängenbereich
	\end{itemize}
	
	\subsubsection{T0-Modell: Universelle Energiefluss-Interaktion}
	\label{subsubsec:t0_universeller_fluss}
	
	Das T0-Modell offenbart Lichtstärke als elektromagnetische Energiefluss-Interaktion mit dem universellen Zeitfeld.
	
	\subsection{T0-Herleitung der Lichtstärke}
	\label{subsec:t0_candela_herleitung}
	
	\subsubsection{Photon-Zeitfeld-Interaktion}
	\label{subsubsec:photon_zeitfeld}
	
	Für elektromagnetische Strahlung wird das T0-Zeitfeld zu:
	\begin{equation}
		T_{\text{photon}}(\vec{x},t) = \frac{1}{\max(E_{\text{photon}}, \omega)}
		\label{eq:photon_zeitfeld}
	\end{equation}
	
	\subsubsection{Visueller Energiebereich im T0-Framework}
	\label{subsubsec:visueller_energiebereich}
	
	Menschliches Sehen operiert im Bereich $\Evis \approx 1.8 - 3.1$ eV. Der T0-Skalierungsparameter für diesen Bereich:
	\begin{equation}
		\xipar_{\text{visuell}} = 2\sqrt{G} \cdot \Evis = 2\sqrt{G} \cdot (2.4 \text{ eV}) \approx 3.9 \times 10^{-28}
		\label{eq:xi_visuell}
	\end{equation}
	
	\subsubsection{T0-Lichtstärke-Formel}
	\label{subsubsec:t0_lichtstaerke_formel}
	
	Die vollständige T0-Herleitung ergibt:
	\begin{equation}
		\boxed{I_{\text{T0}} = \Cto \cdot \frac{\Evis}{\EP} \cdot \Phiphoton \cdot \etavis(\lambda)}
		\label{eq:t0_candela_fundamental}
	\end{equation}
	
	\textbf{Physikalische Komponenten:}
	\begin{itemize}
		\item $\Cto \approx 683$ lm/W: T0-Kopplungskonstante (aus Energieverhältnissen hergeleitet)
		\item $\Evis/\EP$: Visuelle Energie relativ zur Planck-Energie
		\item $\Phiphoton$: Elektromagnetischer Energiefluss
		\item $\etavis(\lambda)$: T0-hergeleitete Effizienzfunktion
	\end{itemize}
	
	\subsection{Dimensionsanalyse und Energienatur}
	\label{subsec:candela_dimensional}
	
	\subsubsection{Vollständige Dimensionsanalyse}
	\label{subsubsec:candela_vollstaendige_dimensional}
	
	\begin{align}
		[I_{\text{T0}}] &= [\Cto] \cdot \frac{[E]}{[E]} \cdot [E T^{-1}] \cdot [1] \\
		&= [\text{lm/W}] \cdot [1] \cdot [E T^{-1}] \cdot [1] \\
		&= [\text{lm/W}] \cdot [E^2] \quad \text{(in natürlichen Einheiten wo } [T] = [E^{-1}])
		\label{eq:candela_dimensionsanalyse}
	\end{align}
	
	\subsubsection{Physikalische Interpretation}
	\label{subsubsec:candela_physikalische_interpretation}
	
	Die Candela repräsentiert:
	\begin{equation}
		\text{Candela} = \Cto \times \text{Energiefluss pro Raumwinkel} \times [1] = [\text{lm/W}] \cdot [E T^{-1}] = [\text{lm/W}] \cdot [E^2]
		\label{eq:candela_interpretation}
	\end{equation}
	
	Die Faktoren $\Evis/\EP$ und $\etavis$ sind dimensionslos, der Raumwinkel ebenfalls; $\Cto = 683$~lm/W ist eine feste Konversionskonstante. Die physikalische Dimension trägt allein der Energiefluss $[E T^{-1}] = [E^2]$. Eine zusätzliche $[E^2]$-\emph{Energieinteraktion} mit dem Detektionssystem und eine $[E^3]$-Dimension folgen daraus nicht; sie gehörten zu demselben $[E^2]$-Bild, das für das Mol nicht gilt (siehe Abschnitt zur Dimensionsanalyse der Stoffmenge).
	
	\subsection{T0-Visuelle-Effizienz-Funktion}
	\label{subsec:t0_visuelle_effizienz}
	
	\subsubsection{Energiebasierte Effizienz-Herleitung}
	\label{subsubsec:energie_effizienz_herleitung}
	
	Die visuelle Effizienzfunktion ergibt sich aus T0-Energieskalierung:
	\begin{equation}
		\etavis(\lambda) = \exp\left(-\frac{(E_{\text{photon}} - E_{\text{vis,peak}})^2}{2\sigma_{\text{T0}}^2}\right)
		\label{eq:t0_visuelle_effizienz}
	\end{equation}
	
	wobei:
	\begin{align}
		E_{\text{vis,peak}} &= 2.4 \text{ eV} \quad \text{(T0-vorhergesagtes Maximum)} \\
		\sigma_{\text{T0}} &= \sqrt{\frac{E_{\text{vis,peak}}}{\EP}} \cdot E_{\text{vis,peak}} \quad \text{(T0-hergeleitete Breite)}
	\end{align}
	Mit $\EP = 1{,}221\times10^{28}$~eV ist $\sigma_{\text{T0}} = \sqrt{2{,}4/1{,}221\times10^{28}} \cdot 2{,}4$~eV $= 3{,}4\times10^{-14}$~eV. Die Funktion ist damit außerhalb von exakt $2{,}4$~eV praktisch null; ihr Maximum liegt bei $hc/2{,}4\,\text{eV} = 517$~nm, nicht bei 555~nm.
	
	\subsubsection{Verbindung zur T0-Kopplungskonstante}
	\label{subsubsec:t0_kopplungskonstante}
	
	Das T0-Modell sagt die Kopplungskonstante vorher:
	\begin{equation}
		\Cto = 683 \text{ lm/W} = f\left(\frac{\Evis}{\EP}, \xipar_{\text{visuell}}\right)
		\label{eq:t0_kopplungsvorhersage}
	\end{equation}

	Dabei bezeichnet $f$ eine allgemeine, nicht ausgeschriebene Abhängigkeit, keine Frequenz und nicht die Grundwindungszahl $f = 1/\xi$.

	Dies liefert eine fundamentale Herleitung des scheinbar willkürlichen 683-lm/W-Faktors.
	
	\section{Universelle Energiebeziehungen: Vollständige Analyse}
	\label{sec:universelle_energiebeziehungen}
	
	\subsection{Alle SI-Einheiten: Energiebasierte Klassifikation}
	\label{subsec:alle_si_energiebasiert}
	
	\subsubsection{Vollständige T0-Abdeckung}
	\label{subsubsec:vollstaendige_t0_abdeckung}
	
	\begin{table}[htbp]
		\resizebox{\textwidth}{!}{%
		\centering
		\begin{tabular}{lcccl}
			\toprule
			\textbf{SI-Einheit} & \textbf{T0-Beziehung} & \textbf{Energie-Dim.} & \textbf{T0-Parameter} & \textbf{Status} \\
			\midrule
			Sekunde (s) & $T = 1/E$ & $[E^{-1}]$ & Direkt & Fundamental \\
			Meter (m) & $L = 1/E$ & $[E^{-1}]$ & Direkt & Fundamental \\
			Kilogramm (kg) & $M = E$ & $[E]$ & Direkt & Fundamental \\
			Kelvin (K) & $\Theta = E$ & $[E]$ & Direkt & Fundamental \\
			Ampere (A) & $I \propto E_{\text{Ladung}}$ & Komplex & $\xipar_{\text{EM}}$ & Elektromagnetisch \\
			\rowcolor{blue!10}
			Mol (mol) & $n = \int \rhoE/\Echar$ & $[1]$ & $\xipar_{\text{atomar}}$ & \textbf{T0-Hergeleitet} \\
			\rowcolor{blue!10}
			Candela (cd) & $I_v \propto \Evis \Phiphoton/\EP$ & $[E^2]$ (Fluss) & $\xipar_{\text{visuell}}$ & \textbf{T0-Hergeleitet} \\
			\bottomrule
		\end{tabular}}
		\caption{Vollständige T0-Modell-Energieabdeckung aller 7 SI-Basiseinheiten}
		\label{tab:vollstaendige_t0_si_abdeckung}
	\end{table}
	
	\subsubsection{Revolutionäre Implikation}
	\label{subsubsec:revolutionaere_implikation}
	
	\begin{tcolorbox}[colback=green!5!white,colframe=green!75!black,title=T0-Modell: Universelles Energieprinzip bestätigt]
		\textbf{Alle 7/7 SI-Basiseinheiten haben fundamentale Energiebeziehungen.}
		
		Es gibt keine \textit{Nicht-Energie}-physikalischen Größen. Die scheinbaren Grenzen waren Artefakte konventioneller Definitionen, nicht fundamentaler Physik.
		
		\textbf{Energie ist die universelle physikalische Größe, aus der alle anderen hervorgehen.}
	\end{tcolorbox}
	
	\subsection{T0-Parameter-Hierarchie}
	\label{subsec:t0_parameter_hierarchie}
	
	\subsubsection{Energieskala-Hierarchie}
	\label{subsubsec:energieskala_hierarchie}
	
	Die T0-Skalierungsparameter umspannen die vollständige Energiehierarchie:
	
	\begin{align}
		\xipar_{\text{Planck}} &= 2\sqrt{G} \cdot \EP = 2 \\
		\xipar_{\text{elektroschwach}} &= 2\sqrt{G} \cdot (100 \text{ GeV}) \approx 1.6 \times 10^{-17} \\
		\xipar_{\text{QCD}} &= 2\sqrt{G} \cdot (1 \text{ GeV}) \approx 1.6 \times 10^{-19} \\
		\xipar_{\text{visuell}} &= 2\sqrt{G} \cdot (2.4 \text{ eV}) \approx 3.9 \times 10^{-28} \\
		\xipar_{\text{atomar}} &= 2\sqrt{G} \cdot (1 \text{ eV}) \approx 10^{-28}
	\end{align}
	
	\subsubsection{Universelle Skalierungsverifikation}
	\label{subsubsec:universelle_skalierungsverifikation}
	
	Das T0-Modell sagt universelle Skalierungsbeziehungen vorher:
	\begin{equation}
		\frac{\xipar(E_1)}{\xipar(E_2)} = \frac{E_1}{E_2}
		\label{eq:universeller_skalierungstest}
	\end{equation}
	
	Mit $G = 1/E_P^2$ ist $\xipar(E) = 2\sqrt{G}\,E = 2E/E_P$; die Skalierung ist daher linear. In SI mit $\sqrt{G}$ in m$^{3/2}$kg$^{-1/2}$s$^{-1}$ und $E$ in Joule ergäbe sich keine reine Zahl; die Zahlenbeispiele unten sind deshalb mit $E_P$ in eV gerechnet.
	
	Dies liefert strenge experimentelle Tests über alle Energieskalen.
	
	\section{T0-Modell-Berechnete Werte}
	\label{sec:t0_berechnete_werte}
	
	\subsection{Mol: Spezielle numerische Ergebnisse}
	\label{subsec:mol_numerische_ergebnisse}
	
	\subsubsection{Standard-Testfall: 1 Mol Wasserstoffatome}
	\label{subsubsec:mol_wasserstoff_test}
	
	\textbf{Eingabeparameter:}
	\begin{itemize}
		\item Charakteristische Energie: $\Echar = 1.0$ eV $= 1.602 \times 10^{-19}$ J
		\item Volumen bei STP: $V = 0.0224$ m³
		\item Avogadro-Zahl: $N_A = 6.022 \times 10^{23}$ mol$^{-1}$
	\end{itemize}
	
	\textbf{T0-Berechnung:}
	\begin{align}
		E_{\text{gesamt}} &= N_A \times \Echar = 6.022 \times 10^{23} \times 1.602 \times 10^{-19} = 9.647 \times 10^{4} \text{ J} \\
		\rhoE &= \frac{E_{\text{gesamt}}}{V} = \frac{9.647 \times 10^{4}}{0.0224} = 4.306 \times 10^{6} \text{ J/m}^3 \\
		n_{\text{T0}} &= \frac{1}{N_A} \int_V \frac{\rhoE}{\Echar} \, d^3x = \frac{1}{N_A} \times \frac{\rhoE \times V}{\Echar} = \frac{4.306 \times 10^{6} \times 0.0224}{1.602 \times 10^{-19}} \times \frac{1}{N_A}
	\end{align}
	
	\textbf{T0-Ergebnis:}
	\begin{equation}
		\boxed{n_{\text{T0}} = 1.000000 \text{ mol (nach SI-Definition von } N_A\text{)}}
		\label{eq:mol_t0_ergebnis}
	\end{equation}
	
	\textbf{T0-Errungenschaft:} Energiebasierte Formulierung der Stoffmenge, keine numerische Vorhersage
	
	\subsubsection{T0-Skalierungsparameter}
	\label{subsubsec:mol_skalierungsparameter}
	
	\begin{equation}
		\xipar_{\text{atomar}} = 2\sqrt{G} \times \Echar = \frac{2 \times 1\,\text{eV}}{1.221 \times 10^{28}\,\text{eV}} = \mathbf{1.64 \times 10^{-28}}
		\label{eq:xi_atomar_berechnet}
	\end{equation}
	
	\subsubsection{Dimensionale Verifikation}
	\label{subsubsec:mol_dimensionale_verifikation}
	
	Mit der Planck-Skala verknüpft ergibt sich die Zahl:
	\begin{equation}
		[n_{\text{T0}}]_{\text{tief}} = \left[\frac{E_{\text{gesamt}}}{\Echar}\right] \times \left[\frac{\Echar}{\EP}\right]^2 = 4.040 \times 10^{-33} \text{ [dimensionslos]}
		\label{eq:mol_e2_dimension}
	\end{equation}
	Das ist $N_A \cdot (1\,\text{eV}/E_P)^2 = 4{,}040\times10^{-33}$. Beide Faktoren sind dimensionslos, das Ergebnis ebenso; eine $[E^2]$-Dimension wird damit nicht belegt (siehe Abschnitt zur Dimensionsanalyse).
	
	\subsection{Candela: Spezielle numerische Ergebnisse}
	\label{subsec:candela_numerische_ergebnisse}
	
	\subsubsection{Standard-Testfall: 1 Watt bei 555 nm}
	\label{subsubsec:candela_555nm_test}
	
	\textbf{Eingabeparameter:}
	\begin{itemize}
		\item Maximale visuelle Wellenlänge: $\lambda = 555$ nm
		\item Photonenenergie: $E_{\text{photon}} = hc/\lambda = 2.234$ eV
		\item Visuelle Energieskala: $\Evis = 2.4$ eV $= 3.845 \times 10^{-19}$ J
		\item Strahlungsfluss: $\Phiphoton = 1.0$ W
	\end{itemize}
	
	\textbf{T0-Berechnung:}
	\begin{align}
		\Cto &= 683 \text{ lm/W} \quad \text{(T0-hergeleitete Kopplungskonstante)} \\
		\frac{\Evis}{\EP} &= \frac{3.845 \times 10^{-19}}{1.956 \times 10^{9}} = 1.966 \times 10^{-28} \\
		\etavis(555\text{nm}) &= 1.0 \quad \text{(SI-Definition } V(555\,\text{nm}) = 1\text{)} \\
		I_{\text{T0}} &= \Cto \times \Phiphoton \times \etavis = 683 \times 1.0 \times 1.0
	\end{align}
	Der Wert $1{,}0$ entspricht der SI-Definition $V(555\,\text{nm}) = 1$, nicht der T0-Formel: Mit der Formel für $\etavis$ oben ergibt sich bei $E_{\text{photon}} = 2{,}234$~eV $\exp\!\left(-0{,}166^2/(2\cdot(3{,}4\times10^{-14})^2)\right) = \exp(-1{,}2\times10^{25}) = 0$ (numerisch).
	
	\textbf{T0-Ergebnis:}
	\begin{equation}
		\boxed{I_{\text{T0}} = 683.0 \text{ lm \quad (SI: 683\,lm/W; Lichtstärke = lm/sr = cd)}}
		\label{eq:candela_t0_ergebnis}
	\end{equation}
	
	\textbf{T0-Errungenschaft:} Rückführung der Lichtstärke auf einen Energiefluss, keine numerische Vorhersage
	
	\subsubsection{T0-Skalierungsparameter}
	\label{subsubsec:candela_skalierungsparameter}
	
	\begin{equation}
		\xipar_{\text{visuell}} = 2\sqrt{G} \times \Evis = \frac{2 \times 2.4\,\text{eV}}{1.221 \times 10^{28}\,\text{eV}} = \mathbf{3.93 \times 10^{-28}}
		\label{eq:xi_visuell_berechnet}
	\end{equation}
	
	\subsubsection{T0-Kopplungs\-konstanten-Herleitung}
	\label{subsubsec:t0_kopplungsherleitung}
	
	Das T0-Modell sagt die Lichtstrom-Wirkungsgrad-Konstante vorher:
	\begin{equation}
		\Cto = 683 \text{ lm/W} = f\left(\xipar_{\text{visuell}}, \frac{\Evis}{\EP}\right)
		\label{eq:t0_kopplungsvorhersage}
	\end{equation}
	
	Dies liefert eine fundamentale Herleitung des scheinbar willkürlichen 683-lm/W-Faktors aus reinen Energieskalierungsbeziehungen.
	
	\subsubsection{Dimensionale Verifikation}
	\label{subsubsec:candela_dimensionale_verifikation}
	
	Mit der Planck-Skala verknüpft:
	\begin{equation}
		[I_{\text{T0}}]_{\text{tief}} = \left[\frac{\Evis}{\EP}\right] \times [\Phiphoton] = 1.966 \times 10^{-28} \times \Phiphoton
		\label{eq:candela_e3_dimension}
	\end{equation}
	Der Faktor $\Evis/\EP = 1{,}966\times10^{-28}$ ist dimensionslos; die Dimension des Ergebnisses ist die des Energieflusses $\Phiphoton$, $[E T^{-1}] = [E^2]$ in natürlichen Einheiten, nicht $[E^3]$.
	
	\section{Experimentelles Verifikationsprotokoll}
	\label{sec:experimentelles_verifikationsprotokoll}
	
	\subsection{Mol-Verifikationsexperimente}
	\label{subsec:mol_verifikation}
	
	\subsubsection{Energiedichte-Messprotokoll}
	\label{subsubsec:mol_energie_protokoll}
	
	\textbf{Experimentelle Schritte:}
	\begin{enumerate}
		\item \textbf{Kalorimetrische Messung:} Bestimmung des Gesamtenergiegehalts $\int \rhoE d^3x$
		\item \textbf{Spektroskopische Analyse:} Messung der charakteristischen Teilchenenergie $\Echar$
		\item \textbf{T0-Berechnung:} Berechnung von $n_{\text{T0}}$ unter Verwendung von \cref{eq:t0_mol_fundamental}
		\item \textbf{Vergleich:} Vergleich mit konventioneller Mol-Bestimmung
	\end{enumerate}
	
	\subsubsection{Vorhergesagte experimentelle Signaturen}
	\label{subsubsec:mol_experimentelle_signaturen}
	
	\begin{itemize}
		\item Energieabhängigkeit: $n_{\text{T0}} \propto E_{\text{gesamt}}/\Echar$
		\item Temperaturabhängigkeit: nur über $E_{\text{gesamt}}(T)/\Echar(T)$; ein Gesetz $n_{\text{T0}} \propto T^2$ folgt aus der dimensionslosen Formulierung nicht
		\item Universelle Verhältnisse: $n_{\text{T0}}(A)/n_{\text{T0}}(B) = E_A/E_B$
	\end{itemize}
	
	\subsection{Candela-Verifikationsexperimente}
	\label{subsec:candela_verifikation}
	
	\subsubsection{Energiefluss-Messprotokoll}
	\label{subsubsec:candela_energie_protokoll}
	
	\textbf{Experimentelle Schritte:}
	\begin{enumerate}
		\item \textbf{Radiometrische Messung:} Bestimmung des elektromagnetischen Energieflusses $\Phiphoton$
		\item \textbf{Spektralanalyse:} Messung der Photonen-Energieverteilung
		\item \textbf{T0-Berechnung:} Anwendung der T0-visuellen Effizienzfunktion \cref{eq:t0_visuelle_effizienz}
		\item \textbf{Intensitätsberechnung:} Berechnung von $I_{\text{T0}}$ unter Verwendung von \cref{eq:t0_candela_fundamental}
		\item \textbf{Vergleich:} Vergleich mit konventioneller Candela-Messung
	\end{enumerate}
	
	\subsubsection{Vorhergesagte experimentelle Signaturen}
	\label{subsubsec:candela_experimentelle_signaturen}
	
	\begin{itemize}
		\item Energiefluss-Abhängigkeit: $I_{\text{T0}} \propto \Phiphoton$
		\item Wellenlängen-Skalierung: $I_{\text{T0}}(\lambda) \propto E_{\text{photon}}(\lambda)$
		\item Universelle Effizienz: $\etavis(\lambda)$ folgt T0-Energieskalierung
	\end{itemize}
	
	\section{Theoretische Implikationen und Vereinheitlichung}
	\label{sec:theoretische_implikationen}
	
	\subsection{Lösung fundamentaler Physikprobleme}
	\label{subsec:loesung_fundamentaler_probleme}
	
	\subsubsection{Das \textit{Nicht-Energie}-Größen-Problem}
	\label{subsubsec:nicht_energie_problem_geloest}
	
	\textbf{Problem gelöst:} Es existieren keine physikalischen Größen ohne Energiebeziehungen.
	
	\textbf{Früheres Missverständnis:} Mol und Candela schienen Ausnahmen von der Energie-Universalität zu sein.
	
	\textbf{T0-Lösung:} Beide Größen haben energiebasierte Herleitungen; die Stoffmenge ist dabei ein dimensionsloses Energieverhältnis.
	
	\subsubsection{Einheitensystem-Vereinheitlichung}
	\label{subsubsec:einheitensystem_vereinheitlichung}
	
	Das T0-Modell liefert die erste wahrhaft vereinheitlichte Beschreibung aller physikalischen Einheiten:
	
	\begin{itemize}
		\item \textbf{Universelle Energiebasis:} Alle 7 SI-Einheiten energiehergeleitet
		\item \textbf{Einzelner Skalierungsparameter:} $\xipar = 2\sqrt{G} \cdot E$
		\item \textbf{Hierarchie-Erklärung:} Verschiedene Energieskalen, dieselbe Physik
		\item \textbf{Experimentelle Einheit:} Universelle Skalierungstests über alle Einheiten
	\end{itemize}
	
	\subsection{Verbindung zur Quantenfeldtheorie}
	\label{subsec:qft_verbindung}
	
	\subsubsection{Teilchenzahl-Operator}
	\label{subsubsec:teilchenzahl_operator}
	
	Die T0-Mol-Herleitung verbindet direkt mit der QFT:
	\begin{equation}
		n_{\text{T0}} \leftrightarrow \langle \hat{N} \rangle = \left\langle \int \hat{\psi}^\dagger(\vec{x}) \hat{\psi}(\vec{x}) d^3x \right\rangle
		\label{eq:mol_qft_verbindung}
	\end{equation}
	
	\subsubsection{Elektromagnetische Feldenergie}
	\label{subsubsec:em_feldenergie}
	
	Die T0-Candela-Herleitung verbindet mit der elektromagnetischen Feldtheorie:
	\begin{equation}
		I_{\text{T0}} \leftrightarrow \mathcal{H}_{\text{EM}} = \frac{1}{2}\int (\vec{E}^2 + \vec{B}^2) d^3x
		\label{eq:candela_em_verbindung}
	\end{equation}
	
	\subsection{Kosmologische und fundamentale Skala-Verbindungen}
	\label{subsec:kosmologische_verbindungen}
	
	\subsubsection{Planck-Skala-Entstehung}
	\label{subsubsec:planck_skala_entstehung}
	
	Sowohl Mol als auch Candela lassen sich mit der Planck-Skala verknüpfen; für das Mol gilt das nur über den Ansatz für $N_A^{(\text{T0})}$ (\cref{eq:avogadro_t0_vorhersage}) mit nicht hergeleitetem $\mathcal{C}_{\text{T0}}$:
	
	\begin{align}
		\text{Mol:} \quad &n_{\text{T0}} = \left(\frac{\Echar}{\EP}\right)^2 \frac{E_{\text{gesamt}}}{\Echar\,\mathcal{C}_{\text{T0}}} \quad \text{(dimensionslos, Ansatz)} \\
		\text{Candela:} \quad &I_{\text{T0}} \propto \frac{\Evis}{\EP} \cdot \Phiphoton
	\end{align}
	
	\subsubsection{Universelle Konstanten aus T0}
	\label{subsubsec:universelle_konstanten_t0}
	
	Das T0-Modell sagt fundamentale Konstanten vorher:
	\begin{align}
		N_A &= f\left(\frac{\Echar}{\EP}\right) \quad \text{(Avogadro-Zahl)} \\
		683 \text{ lm/W} &= g\left(\frac{\Evis}{\EP}\right) \quad \text{(Lichtstrom-Wirkungsgrad)}
	\end{align}
